\documentclass[11pt]{article}
\usepackage[left=0.85in,right=0.85in,top=0.7in,bottom=0.7in]{geometry}
\usepackage[T1]{fontenc}
\usepackage{lmodern}
\usepackage{amsmath,amssymb,amsthm,mathtools,booktabs}
\usepackage[colorlinks=true,linkcolor=blue,citecolor=blue,urlcolor=blue]{hyperref}
\usepackage{microtype}
\newtheorem{theorem}{Theorem}
\newtheorem{lemma}[theorem]{Lemma}
\newcommand{\Q}{\mathbb Q}
\newcommand{\Z}{\mathbb Z}
\newcommand{\HH}[2]{H_{#1}^{(#2)}}
\newcommand{\dd}{\,\mathrm{d}}
\title{Growth comparisons and the arithmetic cost\\of shifting odd-zeta forms}
\author{}
\date{27 October 2026}
\begin{document}
\maketitle
\begin{abstract}
We apply polynomial and logarithmic growth comparisons to the actual
primitive coefficients of tail-moment determinants for odd zeta values.
An elementary prime-denominator argument gives an explicit lower bound
for the primitive target coefficient, including the full gcd reduction.
It proves that shifts satisfying
$N^2\log(a_N+2N)=o(a_N)$ make the normalized height diverge rather than
vanish. Thus increasing the shift sufficiently rapidly cannot supply the
missing irrationality estimate. Lean verifies the growth-comparison and
rational-denominator lemmas described below. The estimate required for
smaller shifts remains unproved, and no all-odd irrationality conclusion
is claimed.
\end{abstract}

\section{The coefficients to be compared}
Fix $s=2r+1$, $r\ge1$, $N\ge1$, and $a\ge0$. Put
\[
 H_k^{(s)}=\sum_{j=1}^k j^{-s},\qquad
 \Delta_{N,a}=\det[\zeta(s)-H_{a+i+j}^{(s)}]_{0\le i,j<N}
             =A_{N,a}\zeta(s)+B_{N,a}.
\]
Subtracting row zero from every other row proves the affine identity.
The gamma integral gives the positive weight
\[
 w_s(x)=\frac{(-\log x)^{s-1}}{(s-1)!(1-x)},\qquad
 \zeta(s)-H_k^{(s)}=\int_0^1x^k w_s(x)\dd x .
\]
In the basis $1,1-x,x(1-x),\ldots,x^{N-2}(1-x)$, whose change-of-basis
determinant is $\pm1$, the target coefficient is the lower Gram determinant
\begin{equation}
 A_{N,a}=
 \det[(a+i+j+1)^{-s}-(a+i+j+2)^{-s}]_{0\le i,j<N-1}.
 \label{eq:A}
\end{equation}
It is positive. Hadamard's inequality and diagonal entries at most one
give $A_{N,a}\le1$; for $N=1$, the empty determinant is one.

Write the rational approximation in lowest terms as
\[
 -\frac{B_{N,a}}{A_{N,a}}=\frac{p_{N,a}}{u_{N,a}},\qquad
 p_{N,a}\in\Z,\quad u_{N,a}\in\Z_{>0}.
\]
Put $Q=\operatorname{lcm}(1,\ldots,a+2N-2)^s$, with an empty lcm equal to one.
Then $U=Q^NA$ and $V=Q^NB$ are integers. For
$g=\gcd(|U|,|V|)$, we have $u=U/g$ and $p=-V/g$. Thus the following bounds
concern the fully reduced coefficient, not an unnormalized presentation.

The endpoint calculation has the exact minimum
\[
 E_{N,a}:=\frac{\Delta_{N,a}}{A_{N,a}}
 =\min_{\deg P<N,\ P(1)=1}\int_0^1P(x)^2x^aw_s(x)\dd x .
\]
Comparing $1-x\le-\log x\le(1-x)/x$ with the Jacobi minimum
\[
 \min_{\deg P<N,\ P(1)=1}\int_0^1P^2x^b(1-x)^m\dd x
 =\frac{m!(m+1)!}{(N)_{m+1}(N+b)_{m+1}}
\]
gives, for $N>r$,
\[
 \frac{(2r-1)!}{(N)_{2r}(N+a)_{2r}}
 \le E_{N,a}\le
 \frac{(2r-1)!}{(N-r)_{2r}(N-r+a)_{2r}},
 \qquad (z)_k=\prod_{j=0}^{k-1}(z+j).
\]
The lower comparison uses $m=2r-1$, $b=a$; the upper one restricts
$P=x^rQ$. Consequently
$N^{2r}(N+a_N)^{2r}E_{N,a_N}\to(2r-1)!$ uniformly for nonnegative shifts.
Small primitive forms therefore require exactly
\begin{equation}
 \frac{u_{N,a_N}}{N^{2r}(N+a_N)^{2r}}\longrightarrow0.
 \label{eq:target}
\end{equation}

\section{What growth comparison proves}
If $a_N\ge0$ and $|u_N|\le C N^p$ eventually, then
\[
 \left|\frac{u_N}{N^{2r}(N+a_N)^{2r}}\right|
 \le C N^{p-4r}\longrightarrow0\qquad(p<4r).
\]
More generally, a bound $|u_N|\le C N^p(\log N)^k$ with real $k$ still
suffices when $p<4r$, since $(\log N)^k=o(N^\epsilon)$ for every
$\epsilon>0$. Lean verifies both conclusions with the numerator bounds
as explicit hypotheses.

A useful general form, also verified in Lean, assumes a lower envelope
$N+a_N\ge cN^\beta$, $c>0$. Then
\[
 N^{2r}(N+a_N)^{2r}\ge c^{2r}N^{2r(1+\beta)},
\]
so $|u_N|\le CN^p(\log N)^k$ suffices when $p<2r(1+\beta)$.
At equality a negative logarithmic exponent still suffices. This also
handles fluctuating shifts that satisfy the same lower envelope.
When $a_N=o(N)$, the denominator is asymptotic to $N^{4r}$.
For a fixed negative lower bound on $a_N$, it is eventually at least
$2^{-2r}N^{4r}$; the stronger inequality with constant one requires
$a_N\ge0$. If $a_N\asymp N^\beta$, $\beta>1$, the denominator instead
has order $N^{2r(1+\beta)}$. Every comparison must evaluate $u_{N,a_N}$
at that same shift. A numerator estimate for $a=0$ is not uniform in $a$.

A saddle-point calculation for a determinant can estimate its absolute
size, but it does not estimate its coefficient gcd. Likewise a recurrence
must first be proved for the reduced sequence before a recurrence
asymptotic theorem can be applied to it. The next calculation exposes
an arithmetic contribution that these substitutions must retain.

\section{An exact centered denominator budget}
Let
\[
 \delta_k=H_{a+k}^{(s)}-H_a^{(s)},\qquad
 D=\prod_{k=1}^{2N-2}(a+k)^s,\qquad Z_{ij}=-D\delta_{i+j}.
\]
An empty product is one. Each $Z_{ij}$ is an integer, because every
denominator in $\delta_{i+j}$ divides $D$. With $J$ the all-ones matrix,
define
\[
 m=\det(Z+DJ)-\det Z=D^NA_{N,a},\qquad c=\det Z.
\]
The centered determinant identity is
\[
 D^N\Delta_{N,a}=m(\zeta(s)-H_a^{(s)})+c.
\]
Since $0<A_{N,a}\le1$, $m$ is a positive integer and $m\le D^N$.
Combining this identity with the reduced rational approximation yields
\begin{equation}
 H_a^{(s)}=\frac pu+\frac cm.
 \label{eq:center}
\end{equation}
If $b$ is the reduced denominator of $H_a^{(s)}$, then
$b\mid um$: multiply \eqref{eq:center} by $um$.
This divisibility is unaffected by any common factor already removed
from $U,V$.

\begin{lemma}
Put
\[
 T_s(a)=\prod_{\substack{p\le a,\ p\ \mathrm{prime}\\a<3p\\p\nmid 2^s+1}}p.
\]
Then
\begin{equation}
 T_s(a)^s\mid b,\qquad T_s(a)^s\mid um,\qquad
 u\ge\frac{T_s(a)^s}{D^N}.
 \label{eq:primecost}
\end{equation}
\end{lemma}
\begin{proof}
Fix a prime in the displayed product. If $a<2p$, the unique summand
whose denominator is divisible by $p$ is $p^{-s}$. All other summands
have denominators prime to $p$. Hence the reduced denominator contains
$p^s$.

If $2p\le a<3p$, isolate the two summands
\[
 \frac1{p^s}+\frac1{(2p)^s}
 =\frac{2^s+1}{2^s p^s}.
\]
Their reduced denominator contains $p^s$, since $p\nmid2^s+1$.
Adding a rational with denominator prime to $p$ cannot remove that
prime power: if $q-R$ has denominator divisible by $p^s$ and
$p\nmid\operatorname{den}(R)$, the divisibility
$\operatorname{den}(q-R)\mid
 \operatorname{den}(q)\operatorname{den}(R)$ forces
$p^s\mid\operatorname{den}(q)$.
Distinct prime powers are coprime, so their product divides $b$.
Now use $b\mid um$, $0<m\le D^N$, and positivity.
\end{proof}

\section{A proved obstruction to rapidly growing shifts}
Write $\vartheta(x)=\sum_{p\le x}\log p$. The exceptional primes
$p\mid2^s+1$ have product at most $2^s+1$. Therefore
\[
 \log T_s(a)\ge\vartheta(a)-\vartheta(a/3)-\log(2^s+1).
\]
The elementary Chebyshev inequalities, also proved in the pinned
mathlib source \cite{Chebyshev}, are
\[
 \vartheta(a)\ge a\log2-\log(a+1)-2\sqrt a\log a,\qquad
 \vartheta(x)\le x\log4\quad(x\ge0).
\]
They give the explicit estimate
\begin{equation}
 \boxed{\;
 \log u_{N,a}\ge s\left[
 \frac{a\log2}{3}-\log(a+1)-2\sqrt a\log a-\log(2^s+1)
 -N(2N-2)\log(a+2N)\right].
 \;}
 \label{eq:lower}
\end{equation}
Here we used $D^N\le(a+2N)^{sN(2N-2)}$.
No prime number theorem is needed.

\begin{theorem}
For a fixed odd weight $s=2r+1>1$, suppose $a_N\ge1$ and
\[
 \frac{N^2\log(a_N+2N)}{a_N}\longrightarrow0.
\]
Then
\[
 \liminf_{N\to\infty}\frac{\log u_{N,a_N}}{a_N}
 \ge\frac{s\log2}{3},\qquad
 \frac{u_{N,a_N}}{N^{2r}(N+a_N)^{2r}}\longrightarrow+\infty.
\]
In particular, this applies to $a_N=\lfloor N^\beta\rfloor$ for
$\beta>2$, and to exponential shifts.
\end{theorem}
\begin{proof}
The hypothesis implies $a_N/N^2\to\infty$, since
$\log(a_N+2N)\ge1$ eventually. In particular $a_N\to\infty$ and
$a_N\ge N^2$ eventually.
Divide \eqref{eq:lower} by $a_N$. The determinant cost tends to zero
by hypothesis; the other errors do so because
$\log a/a\to0$ and $\log a/\sqrt a\to0$.
This proves the lower limit.
Moreover $\log N+\log(N+a_N)=O(\log a_N)=o(a_N)$.
Subtracting the logarithm of the denominator in \eqref{eq:target}
therefore leaves a quantity tending to $+\infty$.
\end{proof}

\section{Exact certificates and formalization boundary}
The audit checks \eqref{eq:center} and \eqref{eq:primecost} by exact
rational arithmetic at $16$ weight/size/shift choices. It also verifies
$28$ larger certificates of
\[
 u_{N,a}\ge 2^b N^{2r}(N+a)^{2r}
\]
directly from $T_s(a)^s\ge2^bD^NN^{2r}(N+a)^{2r}$.
For the cubic shift $a=N^3$:
\begin{center}
\begin{tabular}{rrrr}
\toprule
$s$ & $N=30$ & $N=40$ & $N=60$\\
\midrule
3 & 873 & 34965 & 245818\\
5 & 1442 & 58260 & 409681\\
7 & 2011 & 81556 & 573545\\
9 & 2580 & 104852 & 737408\\
\bottomrule
\end{tabular}
\end{center}
These are exact finite lower certificates, not fitted growth rates.

The Lean growth module verifies power and logarithmic comparison limits
and their squeeze applications. The harmonic modules verify both the
upper-half and upper-third prime-power divisibilities, their prime-product
versions, and the rational denominator budget. The centered determinant
estimate and asymptotic large-shift theorem are paper proofs here.
Neither the polynomial upper bound for the actual
$u_{N,a_N}$ at smaller shifts nor \eqref{eq:target} has been proved.
The all-odd irrationality result remains outside the conclusions.

\begin{thebibliography}{9}
\small
\bibitem{Chebyshev}
mathlib contributors, \emph{Chebyshev functions and estimates},
Lean source at the project's pinned revision, theorems
\texttt{theta\_ge} and \texttt{theta\_le\_log4\_mul\_x}.
\url{https://github.com/leanprover-community/mathlib4/blob/85e3a25e006c35636f0e53b0e9296caca2685bc0/Mathlib/NumberTheory/Chebyshev.lean}.
\end{thebibliography}
\end{document}
